\chapter{Practical Training and Experiential Learning}

\section{Introduction}
Describe the practical component of the training and its importance in reinforcing theoretical concepts.

\section{Hands-on Training Activities}
Provide a detailed account of practical exercises and activities performed. A table may include S.No., Activity/Practical Topic, Tools Used, and Learning Outcome.

\section{Practical Demonstrations}
Describe demonstrations performed by the trainer and explain the technical procedures or methods observed and practiced.

% ============================================================
% EXAMPLE CHAPTER
% ============================================================

\subsection{Programming Examples}

This chapter provides examples of programs developed during
the internship. The programs can be written in different
programming languages such as Python, C, C++, Java, HTML,
SQL, and JavaScript.

% ============================================================
% 1. PYTHON PROGRAM
% ============================================================

\subsection{Python Program}

The following program demonstrates a simple addition of two
numbers using Python.

\subsubsection{Program}

\begin{code}{Python}
# Program to add two numbers

a = 10
b = 20

sum = a + b

print("Sum =", sum)
\end{code}

\subsubsection{Output}

\begin{verbatim}
Sum = 30
\end{verbatim}


% ============================================================
% 2. C PROGRAM
% ============================================================

\subsection{C Program}

The following program demonstrates the addition of two numbers
using the C programming language.

\subsubsection{Program}

\begin{code}{C}
#include <stdio.h>

int main()
{
    int a = 10;
    int b = 20;
    int sum;

    sum = a + b;

    printf("Sum = %d", sum);

    return 0;
}
\end{code}

\subsubsection{Output}

\begin{verbatim}
Sum = 30
\end{verbatim}


% ============================================================
% 3. C++ PROGRAM
% ============================================================

\subsection{C++ Program}

The following program demonstrates the addition of two numbers
using C++.

\subsubsection{Program}

\begin{code}{C++}
#include <iostream>
using namespace std;

int main()
{
    int a = 10;
    int b = 20;

    int sum = a + b;

    cout << "Sum = " << sum;

    return 0;
}
\end{code}

\subsubsection{Output}

\begin{verbatim}
Sum = 30
\end{verbatim}


% ============================================================
% 4. JAVA PROGRAM
% ============================================================

\subsection{Java Program}

The following program demonstrates the addition of two numbers
using Java.

\subsubsection{Program}

\begin{code}{Java}
class Addition
{
    public static void main(String[] args)
    {
        int a = 10;
        int b = 20;

        int sum = a + b;

        System.out.println("Sum = " + sum);
    }
}
\end{code}

\subsubsection{Output}

\begin{verbatim}
Sum = 30
\end{verbatim}


% ============================================================
% 5. HTML PROGRAM
% ============================================================

\subsection{HTML Program}

The following example demonstrates a simple HTML webpage.

\subsubsection{Program}

\begin{code}{HTML}
<!DOCTYPE html>
<html>

<head>
    <title>My Web Page</title>
</head>

<body>

    <h1>Welcome</h1>

    <p>This is my first web page.</p>

</body>

</html>
\end{code}

\subsubsection{Output}

\begin{figure}[H]
    \centering

    \includegraphics[
        width=0.90\textwidth
    ]{figure/upl-logo.png}

    \caption{Code Output}
    \label{fig:code-output}

\end{figure}


% ============================================================
% 6. SQL PROGRAM
% ============================================================

\subsection{SQL Program}

The following SQL query retrieves students whose marks are
greater than 50.

\subsubsection{Program}

\begin{code}{SQL}
SELECT *
FROM students
WHERE marks > 50
ORDER BY marks DESC;
\end{code}



\section{Experiential Learning Outcomes}
Explain how practical sessions improved technical understanding, problem-solving ability, and confidence in using the selected technology.